\documentclass[11pt]{article}

% ---------- packages ----------
\usepackage[utf8]{inputenc}
\usepackage[T1]{fontenc}
\usepackage{lmodern}
\usepackage[a4paper,margin=1in]{geometry}
\usepackage{amsmath,amssymb,amsthm}
\usepackage{xcolor}
\usepackage[most]{tcolorbox}
\usepackage{booktabs}
\usepackage{array}
\usepackage{enumitem}
\usepackage{titlesec}
\usepackage{parskip}
\usepackage{graphicx}
\usepackage[colorlinks=true,linkcolor=blue!55!black,urlcolor=blue!55!black,citecolor=blue!55!black]{hyperref}

% ---------- colors ----------
\definecolor{accent}{HTML}{1F4E79}
\definecolor{accent2}{HTML}{2E7D32}
\definecolor{accent3}{HTML}{8E24AA}
\definecolor{softgray}{HTML}{F2F4F7}
\definecolor{kvcolor}{HTML}{E8F0FE}
\definecolor{warmc}{HTML}{FFF3E0}

% ---------- section styling ----------
\titleformat{\section}{\Large\bfseries\color{accent}}{\thesection}{0.6em}{}
\titleformat{\subsection}{\large\bfseries\color{accent!85!black}}{\thesubsection}{0.6em}{}
\titlespacing*{\section}{0pt}{1.5em}{0.6em}

% ---------- boxes ----------
\newtcolorbox{buildbox}[1][]{
  colback=kvcolor, colframe=accent, boxrule=0.6pt, arc=2pt,
  left=8pt,right=8pt,top=5pt,bottom=5pt,
  fonttitle=\bfseries, title={Builds on / relation to the shortlist}, #1}

\newtcolorbox{theorybox}[1][]{
  colback=green!5, colframe=accent2, boxrule=0.6pt, arc=2pt,
  left=8pt,right=8pt,top=5pt,bottom=5pt,
  fonttitle=\bfseries, title={Theory target}, #1}

\newtcolorbox{expbox}[1][]{
  colback=warmc, colframe=orange!70!black, boxrule=0.6pt, arc=2pt,
  left=8pt,right=8pt,top=5pt,bottom=5pt,
  fonttitle=\bfseries, title={Minimal first experiment}, #1}

\newtcolorbox{riskbox}[1][]{
  colback=red!4, colframe=red!55!black, boxrule=0.6pt, arc=2pt,
  left=8pt,right=8pt,top=5pt,bottom=5pt,
  fonttitle=\bfseries, title={Main risk \& mitigation}, #1}

\newtcolorbox{claimbox}[1][]{
  colback=accent3!8, colframe=accent3, boxrule=0.8pt, arc=2pt,
  left=8pt,right=8pt,top=6pt,bottom=6pt,
  fonttitle=\bfseries, title={Central claim}, #1}

\newcommand{\supp}{\operatorname{supp}}
\newcommand{\R}{\mathbb{R}}
\newcommand{\term}[1]{\textbf{\color{accent!85!black}#1}}

\begin{document}

% ---------------- TITLE ----------------
\begin{titlepage}
\centering
\vspace*{1.8cm}
{\Huge\bfseries\color{accent} A Sparsity-First Research Agenda\par}
\vspace{0.35cm}
{\Large\bfseries\color{accent} for Collaborative Inference\par}
\vspace{0.9cm}
{\large Six projects bridging tensor / low-rank / sparsity methods with\par edge, federated, and multi-agent LLM systems\par}
\vspace{1.6cm}
\begin{tcolorbox}[colback=accent3!8,colframe=accent3,width=0.9\textwidth,arc=3pt]
\centering\bfseries
Thesis: collaborative inference for high-dimensional data should communicate
\emph{intrinsic structure} --- low-rank factors and sparse innovations ---
rather than dense intermediate tensors, and sparsity should be treated as a
\emph{first-class semantic and scheduling resource}, not merely a way to send fewer bits.
\end{tcolorbox}
\vspace{1.4cm}
\begin{tcolorbox}[colback=softgray,colframe=accent,width=0.9\textwidth,arc=3pt]
\centering
Each project is stated with a formal problem, a concrete method, a theoretical
target, an evaluation protocol, its grounding in the reference papers, and its
main risk. A crosswalk (\S\ref{sec:crosswalk}) shows how the earlier idea
shortlist maps onto these six, and \S\ref{sec:prio} recommends a first project.
\end{tcolorbox}
\vfill
{\normalsize Research prospectus \;\textbullet\; \today\par}
\end{titlepage}

\tableofcontents
\newpage

% ================================================================
\section{Executive summary}

Nine reference papers cluster into three lines of work: (i) \textbf{edge / heterogeneous
collaborative LLM inference} (HexGen~\cite{hexgen}, Jupiter~\cite{jupiter}, the FSU
pipelined-inference poster~\cite{poster}); (ii) \textbf{offline-to-online and multi-agent
RL} (BAQ~\cite{baq}, SADQ~\cite{sadq}, DARE~\cite{dare}, FOCI~\cite{foci}); and (iii)
\textbf{federated learning and longitudinal multi-agent LLM systems}
(SAA~\cite{saa}, the coaching paper~\cite{coaching}). A recurring pattern runs through all
of them: a \emph{high-dimensional object is collapsed to a scalar or a dense tensor}, and
the \emph{structure} of that object --- which coordinates changed, what is low-rank, what
is a rare informative event --- is discarded. That discarded structure is exactly where a
tensor / sparsity research program has leverage.

This document develops six projects that each recover and exploit that structure:

\begin{enumerate}[leftmargin=1.6em,itemsep=1pt]
\item \textbf{SupportAlign} --- support-aware semi-asynchronous federated learning.
\item \textbf{SparsePipe} --- co-adaptive sparse collaborative LLM inference.
\item \textbf{LoSpar-CI} --- progressive low-rank-plus-sparse tensor inference \emph{(flagship)}.
\item \textbf{SparseGraph-FOCI} --- learning the dynamic competitor graph.
\item \textbf{Sparse successor-state models} for collaborative-inference control.
\item \textbf{Sparse temporal tensor memory} for private longitudinal agents.
\end{enumerate}

The unifying move is to replace ``compress to save bits'' with ``represent by structure,''
and to let that structure drive \emph{semantics} (which updates matter), \emph{timing}
(when to aggregate), \emph{scheduling} (what to compute and transmit), and \emph{certification}
(when a coarse answer is provably good enough).

% ================================================================
\section{Shared preliminaries and notation}
\label{sec:prelim}

\begin{description}[leftmargin=0pt,style=nextline,itemsep=2pt]
\item[Dimension and sparsity] A model or signal lives in $\R^{p}$ (or a tensor in $\R^{N}$
with $N=\prod_d n_d$). An update/vector is \emph{$s$-sparse} if it has at most $s$ nonzeros;
its \term{support} is $\supp(x)=\{i:x_i\neq0\}$. The governing empirical fact is $s\ll p$.
\item[Low-rank-plus-sparse] Many intermediate tensors decompose as $A=L+S+E$: a low-rank
part $L$ (shared/global structure), a sparse part $S$ (rare informative events), and small
residual noise $E$. This is the robust-PCA / stable principal component pursuit
model~\cite{rpca}, lifted here to the tensor and streaming setting.
\item[Tensor factorizations] $L$ can be encoded by CP, Tucker, or tensor-train (TT)
factors~\cite{kolda,tt}; communication then scales with factor size, not ambient size $N$.
\item[Sketches] CountSketch~\cite{countsketch} and MinHash~\cite{minhash} give $O(\log p)$-
or $O(k)$-size summaries of high-dimensional vectors/supports that approximately preserve
inner products and set overlaps; the Johnson--Lindenstrauss lemma~\cite{jl} preserves
cosine similarity under random projection.
\item[Reference short-names] HexGen~\cite{hexgen}, Jupiter~\cite{jupiter}, Poster~\cite{poster},
FOCI~\cite{foci}, SAA~\cite{saa}, BAQ~\cite{baq}, SADQ~\cite{sadq}, DARE~\cite{dare},
Coaching~\cite{coaching}.
\end{description}

% ================================================================
\section{Crosswalk: earlier shortlist \texorpdfstring{$\rightarrow$}{->} these six}
\label{sec:crosswalk}

The six projects below are refinements of the earlier idea shortlist, sharpened around
\emph{sparsity as the organizing principle}. The mapping:

\begin{center}
\renewcommand{\arraystretch}{1.3}
\small
\begin{tabular}{@{}p{3.4cm}p{3.3cm}p{4.6cm}p{2.3cm}@{}}
\toprule
\textbf{This project} & \textbf{Earlier shortlist item} & \textbf{Reference papers it extends} & \textbf{Cluster} \\
\midrule
1. SupportAlign & Sketched similarity-preserving federated aggregation & SAA~\cite{saa} $+$ DARE posterior~\cite{dare} & Federated \\
2. SparsePipe & Compression-in-the-loop scheduling (activation axis) & Jupiter~\cite{jupiter}, HexGen~\cite{hexgen}, Poster~\cite{poster} & Edge inference \\
3. LoSpar-CI \textbf{(flagship)} & Compression-in-the-loop scheduling (low-rank$+$sparse KV/activation), generalized & Poster~\cite{poster}, Jupiter~\cite{jupiter}, HexGen~\cite{hexgen} & Edge inference / tensor \\
4. SparseGraph-FOCI & Beyond-mean-field low-rank Q $+$ tensor completion of competitor record & FOCI~\cite{foci} & Multi-agent RL \\
5. Sparse successor models & Tensor-structured behavior / successor models & SADQ~\cite{sadq}, BAQ~\cite{baq}, DARE~\cite{dare} & O2O RL \\
6. Sparse temporal memory & Longitudinal coaching tensor completion & Coaching~\cite{coaching} & LLM agents \\
\bottomrule
\end{tabular}
\end{center}

\noindent The earlier ``integrative vision'' (an RL-driven, compression-aware controller)
is not a separate project here; it is the \emph{endpoint} reached by composing projects
2--5, and is stated as the agenda's central claim (\S\ref{sec:prio}).

% ================================================================
\section{SupportAlign: support-aware semi-asynchronous federated learning}
\label{sec:supportalign}

\subsection*{Context and gap}
SAA~\cite{saa} showed that a client's \emph{influence} on federated aggregation is better
predicted by cosine similarity between the current global model and the stale global model
the client trained from than by round-count staleness, and it uses that similarity to set
aggregation weights and buffer timing. But SAA computes cosine similarity on the
\emph{full} $p$-dimensional model.

\begin{claimbox}
In a $p$-dimensional model with $s$-sparse updates, full-model cosine similarity is
dominated by the overwhelming majority of \emph{unchanged} parameters. Two \emph{harmful}
updates can therefore look almost identical --- their agreement is an artifact of the
shared static bulk, not of the parts that actually changed.
\end{claimbox}

\subsection*{Core problem (formal)}
Client $m$ contributes update $\Delta_m\in\R^{p}$ with support $\Omega_m=\supp(\Delta_m)$,
$|\Omega_m|\le s$. Let $\Delta_g$ be the current global update with support $\Omega_g$.
Ordinary cosine $\cos(\Delta_m,\Delta_g)$ is driven by coordinates in the complement of
$\Omega_m\cup\Omega_g$ (where both are $\approx 0$), so it is uninformative about whether
$\Delta_m$ \emph{helps or fights} the global direction on the coordinates that move.

\subsection*{Proposed method}
Represent each update by its support and its values on that support, and measure influence
with three structure-aware components:
\begin{enumerate}[leftmargin=1.5em,itemsep=1pt]
\item \textbf{Support overlap:} $J(\Omega_m,\Omega_g)=\frac{|\Omega_m\cap\Omega_g|}{|\Omega_m\cup\Omega_g|}$ (Jaccard), estimated from MinHash sketches.
\item \textbf{Signed cosine on the intersecting support:}
$\rho_{\cap}=\cos\!\big(\Delta_m|_{\Omega_m\cap\Omega_g},\,\Delta_g|_{\Omega_m\cap\Omega_g}\big)$ --- whether the client agrees \emph{where it matters}.
\item \textbf{Off-support energy:} $\eta=\|\Delta_m|_{\Omega_g^{c}}\|_2\,/\,\|\Delta_m\|_2$ --- how much of the client's change lies \emph{outside} the global update support (novelty / drift).
\end{enumerate}
Communicate CountSketch/MinHash summaries rather than full supports or full vectors, so
uplink cost scales with the sketch size, not $p$. Feed $(J,\rho_{\cap},\eta)$ into a
\textbf{DARE-style posterior}~\cite{dare} that classifies each update as \emph{aligned},
\emph{novel-but-useful}, or \emph{conflicting} (a three-class generalization of DARE's
two-class Gaussian-posterior exchange rule). The class then sets \emph{both}:
\begin{itemize}[leftmargin=1.5em,nosep]
\item \textbf{Aggregation weights} (SAA-style~\cite{saa}): up-weight aligned; damp conflicting; treat novel-but-useful as exploration.
\item \textbf{Aggregation timing}: trigger a buffer flush when the mix of classes destabilizes the global direction (the support-aware analogue of SAA's directional-similarity threshold).
\end{itemize}

\begin{theorybox}
Make convergence / estimation error scale with $s\log p$ (the sparse effective dimension)
plus an additive \emph{sketching-error} term and a \emph{staleness} term --- rather than
with the ambient $p$. Concretely, target a bound of the form
$\mathbb{E}\,[F(w^k)-F(w^\star)]\le \zeta^k(F(w^0)-F(w^\star)) + c_1\,\tfrac{s\log p}{n} + c_2\,\varepsilon_{\text{sketch}} + c_3\,\bar\tau$,
recovering SAA's guarantee as the dense limit $s\!\to\!p$, $\varepsilon_{\text{sketch}}\!\to\!0$.
\end{theorybox}

\begin{expbox}
Sparse linear regression ($p$ up to $10^5$, controllable $s$) with synthetic clients whose
supports overlap by a tunable fraction: show that support-aware influence separates helpful
from harmful updates when full-model cosine cannot, and that CountSketch summaries preserve
the ranking. Then pruned-CNN / sparse-LoRA updates on FEMNIST and CIFAR.
\end{expbox}

\begin{riskbox}
Sketch collisions could corrupt overlap/inner-product estimates for adversarial supports.
Mitigate with multiple independent hashes (median-of-means), and fall back to exact
support exchange for the (few) clients flagged high-influence.
\end{riskbox}

\begin{buildbox}
Sharpens the earlier ``sketched, similarity-preserving federated aggregation'' idea:
same sketching/JL backbone and SAA extension, but the key modification is that
\emph{sparsity determines semantic influence and timing}, not merely transmitted bits ---
and the DARE posterior gives a principled three-way classification.
\end{buildbox}

% ================================================================
\section{SparsePipe: co-adaptive sparse collaborative LLM inference}
\label{sec:sparsepipe}

\subsection*{Context and gap}
Jupiter~\cite{jupiter} and HexGen~\cite{hexgen} schedule \emph{dense} computation over the
pipeline. But token, neuron, attention, and KV importance are highly \emph{non-uniform} and
\emph{time-varying}: they differ by prompt, by layer, and by decoding step. Current work
picks \emph{either} a sparse model \emph{or} a distributed layout; it does not co-design them.

\subsection*{Core problem}
Given a heterogeneous device set with time-varying bandwidth and load, jointly choose, at
each pipeline boundary and decoding step, \emph{what to compute} and \emph{what to transmit}
so as to meet latency/energy/memory budgets at a target accuracy --- treating sparsity as a
\emph{schedulable resource} rather than a fixed model property.

\subsection*{Proposed method}
\begin{enumerate}[leftmargin=1.5em,itemsep=1pt]
\item \textbf{Predict} a hardware-friendly \emph{block-sparse} activation mask at each pipeline boundary (block granularity chosen for real kernel efficiency, not per-element).
\item \textbf{Jointly choose} the layer partition (as in HexGen/Jupiter), the token budget, the active neuron blocks, the KV blocks to keep/transmit, and the encoding format.
\item \textbf{Adapt the wire format:} send dense tensors when sparsity is low; switch to bitmap / block-index formats when sparsity is high (so index overhead never dominates).
\item \textbf{Cost model with real overheads:} include actual sparse-kernel throughput and index/bitmap overhead --- not idealized FLOP savings --- so the optimizer never chooses ``sparsity'' that is slower in practice.
\item \textbf{Online controller:} a lightweight policy re-plans when bandwidth, device load, or \emph{observed} sparsity drifts (see \S\ref{sec:successor} for a principled successor model of that drift).
\item \textbf{Dense fallback:} preserve exact dense compute for coding, mathematics, or low-confidence outputs (mirroring Jupiter's pluggable/auto-selected outline decoding~\cite{jupiter}).
\end{enumerate}
The optimization objective is multi-term: time-to-first-token (TTFT), time-per-output-token
(TPOT), energy, memory, accuracy, and optionally \emph{worst-device} latency (a fairness
term borrowed from FOCI~\cite{foci}).

\begin{theorybox}
An accuracy-vs-sparsity guarantee per boundary: bound the output perturbation induced by a
block-sparse mask via layer Lipschitz constants, and expose it as a constraint the scheduler
must satisfy --- so the plan is chosen subject to a provable accuracy floor rather than a
heuristic threshold.
\end{theorybox}

\begin{expbox}
On the Poster's 3-device Jetson Orin cluster~\cite{poster} (or Jupiter's Jetson setup),
add contextual activation/KV block-sparsity with an adaptive dense/bitmap wire format;
measure TPOT and quality (GPT-judge on WizardLM/Vicuna) versus the dense pipeline at matched
memory, and show the online controller recovers when one link is throttled mid-generation.
\end{expbox}

\begin{riskbox}
Realized speedup from unstructured sparsity is often eaten by kernel and index overhead.
Mitigate by committing to \emph{block} sparsity sized to the hardware and by putting measured
kernel efficiency \emph{inside} the cost model, so unprofitable sparsity is never selected.
\end{riskbox}

\begin{buildbox}
This is the ``activation axis'' of the earlier compression-in-the-loop scheduling idea:
the distinctive contribution is treating sparsity as a \emph{time-varying scheduling
resource} and co-designing it with the distributed layout, which neither Jupiter, HexGen,
nor the Poster does.
\end{buildbox}

% ================================================================
\section{LoSpar-CI: progressive low-rank-plus-sparse tensor inference (flagship)}
\label{sec:lospar}

This is the project most aligned with a tensor-research identity, and the one with the best
chance of founding a recognizable line rather than another scheduling variant.

\subsection*{Core assumption}
An intermediate high-dimensional tensor at layer/stage $\ell$ decomposes as
\[
A_\ell \;=\; L_\ell \;+\; S_\ell \;+\; E_\ell,
\]
where $L_\ell$ is low-rank shared structure, $S_\ell$ contains rare informative events
(sparse), and $E_\ell$ is small residual noise.

\subsection*{Proposed system}
\begin{enumerate}[leftmargin=1.5em,itemsep=1pt]
\item \textbf{Encode} $L_\ell$ with CP / Tucker / tensor-train factors~\cite{kolda,tt}; encode $S_\ell$ with support indices $+$ values.
\item \textbf{Send coarse first:} transmit a coarse low-rank approximation of $A_\ell$ across the collaborating devices.
\item \textbf{Compute-and-certify:} run the downstream prediction on the coarse tensor and compute an uncertainty / output-margin bound.
\item \textbf{Refine on demand:} request additional sparse residual blocks of $S_\ell$ \emph{only if} the coarse prediction is not certifiably stable.
\item \textbf{Joint selection:} choose rank, sparse budget, pipeline cut, and target device together (the LoSpar analogue of HexGen's placement search~\cite{hexgen}).
\end{enumerate}

\begin{theorybox}
If the downstream map is Lipschitz and the output margin exceeds a function of the
reconstruction error $\|A_\ell-(L_\ell+S_\ell)\|$, then the coarse prediction is
\textbf{certifiably identical} to full dense inference --- no accuracy loss, only earlier
stopping. Communication then scales as (factor size) $+\;k\log(N/k)$ for a $k$-sparse
residual, rather than the ambient tensor size $N$.
\end{theorybox}

\subsection*{Best application domains}
Ultra-high-dimensional signals where the $A=L+S+E$ model is physically justified:
whole-slide pathology images; hyperspectral imagery; 3D medical volumes; multiview sensor
networks; long video / scientific-simulation tensors.

\subsection*{Positioning}
Token-compression methods already shrink large-image inputs --- e.g. DTC-WSI~\cite{dtcwsi}
for whole-slide images and SparseVLM~\cite{sparsevlm} for VLM token pruning --- but those
prune tokens for a \emph{single} model. \emph{Progressive low-rank-plus-sparse transmission
across heterogeneous collaborating devices, with a certified early-stop rule}, is
materially different: the object communicated is intrinsic structure, and refinement is
demand-driven and provably safe.

\begin{claimbox}
Collaborative inference for ultra-high-dimensional data should communicate intrinsic
structure --- low-rank factors and sparse innovations --- rather than dense intermediate
tensors.
\end{claimbox}

\begin{expbox}
Pick \emph{one} modality (whole-slide images or hyperspectral), \emph{one} decomposition
(e.g. tensor-train), \emph{one} progressive refinement rule (margin-based early stop), and a
\emph{two- or three-device} prototype. Show communication reduction at zero (certified) or
bounded accuracy loss versus dense transmission; report the fraction of inputs that stop at
the coarse stage. Defer the scheduling/RL extension until this structural premise holds.
\end{expbox}

\begin{riskbox}
The $A=L+S+E$ premise may hold weakly for some layers/modalities, weakening certified
early-stop. Mitigate by measuring per-layer low-rank$+$sparse energy \emph{first} (a cheap
empirical study) and applying LoSpar only where the spectrum justifies it, with dense
fallback elsewhere.
\end{riskbox}

\begin{buildbox}
Generalizes the earlier compression-in-the-loop idea from KV/activation compression to
\emph{arbitrary high-dimensional intermediate tensors}, and adds the decisive new element:
a \emph{certified, progressive} refinement rule rather than fixed-rate compression.
\end{buildbox}

% ================================================================
\section{SparseGraph-FOCI: learning the dynamic competitor graph}
\label{sec:sparsegraph}

\subsection*{Context and gap}
FOCI~\cite{foci} identifies potential competitors via a \emph{geographic perception radius}.
In reality, only a sparse subset of nearby devices meaningfully competes for the same
edge-resource pool, and physical proximity is a noisy proxy under mobility.

\subsection*{Proposed method}
\begin{enumerate}[leftmargin=1.5em,itemsep=1pt]
\item Model device interactions as a \textbf{dynamic sparse hypergraph}: devices are nodes; shared edge-resource pools are hyperedges.
\item Infer the top-$s$ influential competitors of each device from anonymous resource-demand and action \emph{sketches} (privacy preserved, as in FOCI).
\item Impose \textbf{temporal group sparsity} so the graph changes gradually unless evidence supports a structural change (a fused-lasso / group-lasso prior over time).
\item Feed the learned sparse influence distribution into \textbf{mean-field MARL} in place of the geographic neighborhood.
\item Retain FOCI's \textbf{max--min fairness oracle} for tail-QoS control.
\end{enumerate}
An optional extension recovers the earlier ``beyond mean-field'' idea: represent the joint
action-value as a low-rank tensor over agents (mean-field $=$ the rank-1 term) and add a
few low-rank components to capture pairwise contention; use low-rank tensor completion on
the mobility-incomplete competitor record to denoise the statistic.

\begin{theorybox}
Replace FOCI's dependence on the full action-space dimension $\varphi$ (or total device
count) with dependence on the \emph{effective degree} $s$, plus a graph-estimation-error
term and a mobility-rate term --- i.e. re-derive the Serfling-style bound with $\lambda$
instantiated as the learned sparse degree rather than a hand-set radius.
\end{theorybox}

\begin{riskbox}
\textbf{Circularity:} actions determine contention, but the learned contention graph also
determines actions. Mitigate with two-timescale learning (graph updated on a slower
timescale than policies) or delayed graph updates, so the estimator sees a quasi-stationary
action distribution.
\end{riskbox}

\begin{buildbox}
Refines the earlier ``beyond-mean-field / tensor-completion-of-competitor-record'' idea into
a sharper, sparsity-native form: the object learned is a \emph{sparse dynamic hypergraph},
and the payoff is a bound governed by effective degree $s$ instead of ambient dimension.
\end{buildbox}

\begin{expbox}
On FOCI's Porto-taxi $+$ Jetson-profile simulation~\cite{foci}: replace the perception
radius with a learned top-$s$ competitor graph; show tighter worst-case latency/energy tails
and better convergence, ablating the temporal-sparsity prior and the two-timescale schedule.
\end{expbox}

% ================================================================
\section{Sparse successor-state models for collaborative-inference control}
\label{sec:successor}

\subsection*{Context and gap}
SADQ~\cite{sadq} builds Q-targets from a \emph{diagonal-Gaussian} successor-state model.
That model does not reflect high-dimensional \emph{system} states in which only a few
components change between steps --- a device joins/leaves, a queue spikes, a link degrades,
the active model changes. These are \emph{sparse, structured} transitions, not
independent Gaussian jitter.

\subsection*{Proposed transition model}
Factor the successor distribution into \emph{which coordinates change} and \emph{their values}:
\[
P(s'\mid s,a)\;=\;\underbrace{P(\text{support change}\mid s,a)}_{\text{sparse Bernoulli / point process}}\;\cdot\;\underbrace{P(\text{values on support}\mid \text{support},s,a)}_{\text{conditional value model}} .
\]
The support-change factor can be a sparse Bernoulli model, a point process, or a structured
diffusion model; the value factor predicts magnitudes only on the changed coordinates/fibers.
For action selection, use \emph{pessimistic / quantile} aggregation over several predicted
successors (rather than SADQ's max over imagined states), to avoid optimism on rare events.
Train from system logs first, then relax conservatism online with BAQ/DARE-style behavior
alignment~\cite{baq,dare}.

\subsection*{Application and evaluation}
Use the model to choose pipeline cuts, precision, sparse budgets (feeding SparsePipe,
\S\ref{sec:sparsepipe}), device membership, and fallback policies. Evaluate on the metrics
that matter for reliability, not average reward: \emph{tail} SLO violations, adaptation
speed after failures, switching cost, and regret.

\begin{theorybox}
A variance-reduction result in the spirit of SADQ's Theorems 1--2~\cite{sadq}, but for the
factored model: separating support prediction from value prediction lowers target variance
on sparse transitions, and quantile aggregation controls the bias from rare high-impact
events.
\end{theorybox}

\begin{riskbox}
A learned support-change model can be unreliable in fast-changing or under-logged regimes
(SADQ's own failure mode on short-horizon tasks). Mitigate by blending predicted successors
with real ones via a learned confidence weight, defaulting to real transitions when the
model is uncertain.
\end{riskbox}

\begin{buildbox}
Specializes the earlier ``tensor-structured behavior / successor models'' idea to
collaborative-inference control, with the concrete new form being the \emph{support-then-value}
factorization of the transition kernel.
\end{buildbox}

% ================================================================
\section{Sparse temporal tensor memory for private longitudinal agents}
\label{sec:memory}

\subsection*{Context and gap}
The coaching paper~\cite{coaching} externalizes conversation state to fight ``lost in the
middle,'' but treats each user in isolation. Longitudinal user state is naturally a
\emph{sparse temporal tensor}
\[
\text{user}\;\times\;\text{goal}\;\times\;\text{barrier}\;\times\;\text{context}\;\times\;\text{time},
\]
most of whose entries are unchanged at any given turn.

\subsection*{Proposed system}
Maintain a structured, privacy-aware memory:
\begin{itemize}[leftmargin=1.5em,nosep]
\item a \textbf{low-rank component} for stable preferences and recurring patterns (shared archetypes);
\item \textbf{sparse deltas} for new commitments, barriers, or health events (change detection);
\item \textbf{device-local sensitive factors} that never leave the device;
\item \textbf{edge-side anonymized shared factors} for cross-user structure;
\item \textbf{query-dependent fiber retrieval}, so only the relevant slice of state reaches the LLM (bounded context, auditable coverage).
\end{itemize}
This connects tensor completion, sparse change detection, privacy, and collaborative LLM
inference in one object: irregular engagement (missed sessions) is exactly the missing-data
regime tensor completion handles.

\begin{theorybox}
Completion/recovery guarantees for the low-rank $+$ sparse temporal tensor under
\emph{missing-not-at-random} engagement, bounding reconstruction error by observed-fiber
count and archetype rank --- and a privacy statement for the device-local / edge-shared split.
\end{theorybox}

\begin{expbox}
On multi-session coaching logs (or a simulated longitudinal cohort): measure dialogue
continuity, missed-constraint rate, token cost, reconstruction error under withheld
sessions, and privacy leakage of the shared factors --- versus the transcript-append and
memory-only baselines from~\cite{coaching}.
\end{expbox}

\begin{riskbox}
Low-rank archetypes risk over-generalizing across users (stereotyping) and could leak
membership. Mitigate with per-user sparse deltas dominating the personalization signal and
with differential-privacy noise on the shared factors.
\end{riskbox}

\begin{buildbox}
Expands the earlier ``coaching tensor-completion'' idea with the low-rank$+$sparse split,
on-device/edge privacy separation, and query-conditioned fiber retrieval.
\end{buildbox}

% ================================================================
\section{Prioritization and recommended first project}
\label{sec:prio}

\begin{center}
\renewcommand{\arraystretch}{1.3}
\small
\begin{tabular}{@{}p{3.2cm}p{2.4cm}p{2.2cm}p{5.0cm}@{}}
\toprule
\textbf{Project} & \textbf{Strength} & \textbf{Cost} & \textbf{Best fit} \\
\midrule
SupportAlign & Theory-first & Modest & Cleanest math; sharpest ``sparsity $=$ semantics'' story \\
SparsePipe & Systems/LLM & Higher (kernels) & Strongest systems paper; real speedups on edge clusters \\
LoSpar-CI & Distinctive tensor & Medium & Best chance to found a recognizable research line \\
SparseGraph-FOCI & MARL theory$+$sim & Medium & Effective-degree bound; needs stable MARL sim \\
Sparse successor & RL reliability & Medium & Tail-SLO control; depends on SparsePipe as consumer \\
Sparse temporal memory & Cross-cutting & Higher (data/privacy) & Broadest, but gated by longitudinal data access \\
\bottomrule
\end{tabular}
\end{center}

\noindent\textbf{Recommendations.}
\begin{itemize}[leftmargin=1.5em,nosep]
\item \textbf{Best theory-first project:} \emph{SupportAlign} --- a clean mathematical problem, modest implementation, and an unambiguous distinction from ordinary gradient compression.
\item \textbf{Best systems / LLM project:} \emph{SparsePipe} --- directly extends Jupiter/HexGen/the Poster with measurable speedups.
\item \textbf{Best distinctive tensor-and-sparsity project:} \emph{LoSpar-CI} --- the one most likely to establish a recognizable line rather than another scheduling variant.
\end{itemize}

\begin{claimbox}
LoSpar-CI carries the agenda's sharp thesis: \emph{collaborative inference for
ultra-high-dimensional data should communicate intrinsic structure --- low-rank factors and
sparse innovations --- rather than dense intermediate tensors.} A manageable first paper
fixes one modality (e.g. whole-slide or hyperspectral), one decomposition, one progressive
refinement rule, and a two-/three-device prototype; the scheduling/RL extensions
(SparsePipe, Sparse successor) follow once the structural premise is empirically established.
\end{claimbox}

% ================================================================
\begin{thebibliography}{99}
\small
\bibitem{hexgen} Y.~Jiang et al. \emph{HexGen: Generative Inference of LLM over Heterogeneous Environment.} ICML 2024; arXiv:2311.11514.
\bibitem{jupiter} S.~Ye et al. \emph{Jupiter: Fast and Resource-Efficient Collaborative Inference of Generative LLMs on Edge Devices.} IEEE INFOCOM 2025; arXiv:2504.08242.
\bibitem{poster} H.~Zhou and X.~Zhang. \emph{Poster: Distributed Pipelined Collaborative Inference for LLMs on Edge Device Clusters.} ACM HotMobile 2026.
\bibitem{foci} H.~Zhou and X.~Zhang. \emph{Fairness-Oracular MARL with Competitor-Aware Signals for Collaborative Inference.} NeurIPS 2025 Workshop (AI4NextG).
\bibitem{saa} H.~Zhou, J.~Fu, Y.~Yuan, L.~Guo, X.~Zhang. \emph{Similarity-Guided Rapid Deployment of Federated Intelligence over Heterogeneous Edge Computing.} IEEE INFOCOM 2025.
\bibitem{baq} L.~Zu, H.~Zhou, X.~Zhang. \emph{Behavior-Adaptive Q-Learning: A Unifying Framework for Offline-to-Online RL.} arXiv:2511.03695, 2025.
\bibitem{sadq} L.~Zu, H.~Zhou, X.~Zhang. \emph{Enhancing Q-Value Updates in Deep Q-Learning via Successor-State Prediction.} arXiv:2511.03836, 2025.
\bibitem{dare} L.~Zu, Y.~Qian, S.~Chakraborty, X.~Zhang. \emph{From Static Constraints to Dynamic Adaptation: Sample-Level Constraint Relaxation for Offline-to-Online RL (DARE).} ICML 2026; arXiv:2511.03828.
\bibitem{coaching} H.~Zhou et al. \emph{LLM-based Multi-Agent Coaching for Longitudinal Behavior Change: Multi-Scale Summaries and Dialogue Steering.} Smart Health, 2026.
\bibitem{rpca} E.~J.~Cand\`es, X.~Li, Y.~Ma, J.~Wright. \emph{Robust Principal Component Analysis?} J.~ACM, 2011.
\bibitem{kolda} T.~G.~Kolda and B.~W.~Bader. \emph{Tensor Decompositions and Applications.} SIAM Review, 2009.
\bibitem{tt} I.~V.~Oseledets. \emph{Tensor-Train Decomposition.} SIAM J.~Sci.~Comput., 2011.
\bibitem{countsketch} M.~Charikar, K.~Chen, M.~Farach-Colton. \emph{Finding Frequent Items in Data Streams (CountSketch).} ICALP 2002.
\bibitem{minhash} A.~Z.~Broder. \emph{On the Resemblance and Containment of Documents (MinHash).} SEQUENCES 1997.
\bibitem{jl} W.~B.~Johnson and J.~Lindenstrauss. \emph{Extensions of Lipschitz Mappings into a Hilbert Space.} 1984.
\bibitem{dtcwsi} \emph{DTC-WSI: dynamic token compression for whole-slide-image analysis.} (representative token-compression work for gigapixel pathology).
\bibitem{sparsevlm} \emph{SparseVLM: visual token sparsification for efficient vision-language-model inference.}
\end{thebibliography}

\end{document}
